% Policy Indistinguishability section body.
\subsection{Policy Indistinguishability (PI)}
\label{sec:eval-pi}

PI evaluates whether an observer can identify which of two semantically related policies governs a given execution.
We evaluate PI using $30$ designed policy pairs, divided into three types with $10$ pairs each:
\begin{itemize}
	\item \texttt{semantic\_equivalent}: The two policies share the same logical form, consisting of the same DUO terms and Boolean structure, but differ in predicate parameter values, such as whitelist codes.
	\item \texttt{weak\_equivalent}: The two policies share the same set of DUO terms but differ in Boolean structure.
	\item \texttt{controlled\_difference}: The two policies differ in their sets of DUO terms. For each $k \in \{0,1,2,3,4\}$, we include two pairs whose policies differ in exactly $k$ of the four SPP dimensions (Section~\ref{sec:eval-spp}).
\end{itemize}
Each policy belongs to one pair, and which of its two policies is $P_0$ and which is $P_1$ is drawn at random.

The attacker knows both policies in a pair and can execute each policy to obtain labeled training examples.
Given the on-chain observables of an execution, the attacker predicts whether it was governed by $P_0$ or $P_1$.
For each random seed and policy pair, we split the executions of each policy equally between training and testing, separately for granted and denied outcomes.
We train LR and RF on the training executions and evaluate each classifier using test AUC.
Each execution commits a separate policy instance (Section~\ref{sec:eval-dataset}), so training and test executions never share a commitment.

For each pair type and classifier, we first average the test AUCs across the $10$ pairs for each seed.
To assess significance, we construct a null distribution of this average by repeatedly permuting the test labels within each pair, recomputing each pair's AUC, and averaging across pairs.
This yields one mean AUC and one permutation $p$-value per seed.
We then apply the leakage criterion (Section~\ref{sec:eval-method}) using the AUC averaged over the $30$ seeds and the fraction of seeds with $p<0.05$.

All three pair types are classified as PASS (Table~\ref{tab:pi-pairs}).
The highest mean AUC, $0.549$ for \texttt{semantic\_equivalent} pairs with RF, remains below the $0.55$ threshold, with $20\%$ of seeds significant.
